\clearpage\section{MOS transistor schematic}
These editable vector schematics connect equal node names within a circuit. Primitive-cell names are local pins; the hierarchy defines their connections. B denotes the body terminal; 0 is global ground. Each panel identifies its netlist and scope; measurement sources are connected as described.
The MOS primitives and full dff hierarchy in \texttt{02\_setup\_sweep.sp} and \texttt{03\_hold\_sweep.sp}.
\par
\begin{figure}[!ht]
\centering
\begin{circuitikz}[font=\small]
\node[pmos,bulk] (MP) at (0,2.5) {};
\node[anchor=west,font=\scriptsize] at ($(MP.center)+(.55,.40)$) {MP};
\draw (MP.bulk) -- ++(.30,0) node[right,font=\scriptsize] {B=vp};
\node[nmos,bulk] (MN) at (0,0.2) {};
\node[anchor=west,font=\scriptsize] at ($(MN.center)+(.55,.40)$) {MN};
\draw (MN.bulk) -- ++(.30,0) node[right,font=\scriptsize] {B=vn};
\draw (MP.G) -- (-1,2.5) -- (-1,.2) -- (MN.G); \draw (-1,1.35) node[circ]{} -- (-1.6,1.35) node[left]{a}; \draw (MP.D |- 0,1.35) node[circ]{} -- (2,1.35) node[right]{z};
\draw (MP.S) -- ++(0,0.5) node[above,font=\small] {\texttt{vp}};
\draw (MP.D) -- (MN.D) ;
\draw (MN.S) -- ++(0,-0.5) node[below,font=\small] {\texttt{vn}};
\end{circuitikz}
\caption{Exact inv cell; k=1 in the flip-flop.}
\end{figure}
\par
\par
\begin{figure}[!ht]
\centering
\begin{circuitikz}[font=\small]
\node[pmos,bulk,rotate=-90] (TP) at (0,1) {};
\node[nmos,bulk,rotate=90] (TN) at (0,-1) {};
\draw (TP.D) -- (TP.D -| -2,0) -- (TN.D -| -2,0) -- (TN.D);
\draw (TP.S) -- (TP.S -| 2,0) -- (TN.S -| 2,0) -- (TN.S);
\draw (-2,0) node[circ]{} -- (-2.6,0) node[left]{a (D)};
\draw (2,0) node[circ]{} -- (2.6,0) node[right]{z (S)};
\draw (TP.G) -- ++(0,.5) node[above]{eb};
\draw (TN.G) -- ++(0,-.5) node[below]{en};
\draw (TP.bulk) -- ++(0,-.35) -- ++(.8,0) node[right,font=\scriptsize]{B=vp};
\draw (TN.bulk) -- ++(0,.35) -- ++(.8,0) node[right,font=\scriptsize]{B=vn};
\node[font=\scriptsize] at (-.8,1.4){TP};
\node[font=\scriptsize] at (-.8,-1.4){TN};
\end{circuitikz}
\caption{Exact two-MOS transmission gate with all four terminals identified.}
\end{figure}
\par
The inv cell has $W_P=3.6k\,\mu$m, $W_N=1.8k\,\mu$m and $L=0.25\,\mu$m. The tg cell uses the same widths/length. For tg, both device drains are a and both sources are z; the pass path is bidirectional. TP/TN correspond to MP/MN inside the tg subcircuit. PMOS bodies remain vp and NMOS bodies vn, independent of the signal terminals. The n-channel gate is en, and the p-channel gate is eb. Equal node labels connect the two pass devices in parallel.
\clearpage
\par
\begin{figure}[!ht]
\centering
\begin{circuitikz}[font=\small]
\node[draw,rounded corners=2pt,align=center,minimum width=1.7cm,minimum height=.9cm] (X1) at (0,0) {X1: TG\\en=cb};
\node[draw,rounded corners=2pt,align=center,minimum width=1.7cm,minimum height=.9cm] (X2) at (2.6,0) {X2: inv};
\node[draw,rounded corners=2pt,align=center,minimum width=1.7cm,minimum height=.9cm] (X3) at (2.6,-2) {X3: inv};
\node[draw,rounded corners=2pt,align=center,minimum width=1.7cm,minimum height=.9cm] (X4) at (0,-2) {X4: TG\\en=clk};
\node[draw,rounded corners=2pt,align=center,minimum width=1.7cm,minimum height=.9cm] (X5) at (6.0,0) {X5: TG\\en=clk};
\node[draw,rounded corners=2pt,align=center,minimum width=1.7cm,minimum height=.9cm] (X6) at (8.6,0) {X6: inv};
\node[draw,rounded corners=2pt,align=center,minimum width=1.7cm,minimum height=.9cm] (X7) at (8.6,-2) {X7: inv};
\node[draw,rounded corners=2pt,align=center,minimum width=1.7cm,minimum height=.9cm] (X8) at (6,-2) {X8: TG\\en=cb};
\draw (-1.7,0) node[left]{d} -- (X1.west);
\draw[-{Latex}] (X1.east) -- (X2.west) node[midway,above]{m};
\draw[-{Latex}] (X2.east) -- (X5.west) node[midway,above]{mx};
\draw[-{Latex}] (X5.east) -- (X6.west) node[midway,above]{s};
\draw (X6.east) -- (10.4,0) node[right]{q};
\draw[-{Latex}] (3.8,0) node[circ]{} |- (4,-1) -| (X3.east); \draw[-{Latex}] (X3.west) -- (X4.east) node[midway,above]{mf}; \draw[-{Latex}] (X4.west) -- (-1.3,-2) -- (-1.3,-1) -- (1.3,-1) -- (1.3,0) node[circ]{};
\draw[-{Latex}] (9.8,0) node[circ]{} |- (10,-1) -| (X7.east); \draw[-{Latex}] (X7.west) -- (X8.east) node[midway,above]{sf}; \draw[-{Latex}] (X8.west) -- (4.7,-2) -- (4.7,-1) -- (7.3,-1) -- (7.3,0) node[circ]{};
\end{circuitikz}
\caption{Exact master/slave storage connections. Lower feedback inverters carry signal from right to left.}
\end{figure}
\par
Every block uses vp/vn rails; CQ=10 fF connects q to vn. cb is the ideal complement of clk. For each TG, eb is the complement of its shown en: X1/X8 have (en,eb)=(cb,clk), X4/X5 have (clk,cb). X1/X4 form the master input/feedback switches and X5/X8 the slave switches. The feedback path is mx $\to$ X3 $\to$ mf $\to$ X4 $\to$ m; the slave path is q $\to$ X7 $\to$ sf $\to$ X8 $\to$ s. There are four inv and four tg cells, hence 16 MOS. The separate inverter-chain experiment repeats inv four times with k=1 and a 10 fF load; its edge stimulus is specified in its own netlist.
